% Benötigte Pakete:
% \usepackage{longtable}
% \usepackage{booktabs}
% \usepackage{pdflscape}
% \usepackage{array}

\begingroup
\tiny
\renewcommand{\arraystretch}{1}

\setlength{\LTleft}{-1.5cm}
\setlength{\LTright}{0pt}

\begin{longtable}{p{1cm} p{2.5cm} p{3.2cm} p{4.5cm} p{4.5cm}}
\label{tab:kodierleitfaden}\\

\toprule
\textbf{Code} & \textbf{Kategorie} & \textbf{Definition} & \textbf{Codierregel} & \textbf{Ankerbeispiel / Modellbezug} \\
\midrule
\endfirsthead

\toprule
\textbf{Code} & \textbf{Kategorie} & \textbf{Definition} & \textbf{Codierregel} & \textbf{Ankerbeispiel / Modellbezug} \\
\midrule
\endhead

\midrule
\multicolumn{5}{r}{\textit{Fortsetzung auf der nächsten Seite}} \\
\midrule
\endfoot

\bottomrule
\endlastfoot

\multicolumn{5}{l}{\textbf{K1 Qualitätsverständnis und ideale Metadaten}} \\
\midrule

K1.1 & Mindestvalidität und Best Practice
& Aussagen, die zwischen Mindestanforderungen und höherer Qualitätsstufe unterscheiden.
& Codieren, wenn DCAT-AP.de-Validität, Pflichtanforderungen oder Best-Practice-Empfehlungen als Qualitätsstufen beschrieben werden.
& Valide DCAT-AP.de-Lieferung als Minimalvoraussetzung; Best Practice als höhere Qualitätsstufe. Modellbezug: Trennung von Basiskonformität und Zusatzqualität. \\

K1.2 & Qualität statt Quantität
& Aussagen, die betonen, dass nicht möglichst viele, sondern die relevanten Felder in guter Qualität ausgefüllt sein sollten.
& Codieren, wenn ein pragmatisches Qualitätsverständnis beschrieben wird, das relevante und sorgfältig gepflegte Felder höher bewertet als maximale Feldanzahl.
& „Lieber Qualität statt Quantität.“ Modellbezug: keine reine Feldzählung als Qualitätsmaß. \\

K1.3 & Metadaten als Inhalts- und Nutzungsvorschau
& Aussagen, nach denen Metadaten Nutzenden vor dem Öffnen des Datensatzes eine fachliche und technische Einschätzung ermöglichen sollen.
& Codieren, wenn Metadaten als Vorschau auf Inhalt, technische Verarbeitbarkeit oder Nutzungseignung beschrieben werden.
& Metadaten sollen zeigen, was die Daten enthalten und ob sie technisch oder inhaltlich nutzbar sind. Modellbezug: Nutzendenorientierung als Qualitätsanforderung. \\

K1.4 & Semantischer Graph als Idealform
& Aussagen, die einen semantischen Graphen oder Linked-Data-Strukturen als Ideal hochwertiger Metadaten beschreiben.
& Codieren, wenn Graphstruktur, Erweiterbarkeit, Entitäten oder semantische Modellierung als Ideal hochwertiger Metadaten genannt werden.
& Ein perfekter Metadatensatz basiert auf DCAT-AP.de beziehungsweise einem semantischen Graphen. Modellbezug: Graphfähigkeit als übergeordnete Qualitätsanforderung. \\

\midrule
\multicolumn{5}{l}{\textbf{K2 Vollständigkeit, Reichhaltigkeit und Feldrelevanz}} \\
\midrule

K2.1 & Relative Vollständigkeit
& Aussagen, die Vollständigkeit nicht absolut, sondern relativ zum fachlichen Kontext und zu grundlegenden Anforderungen verstehen.
& Codieren, wenn gefragt wird, ob grundlegende Anforderungen und thematisch relevante Informationen vorhanden sind.
& Metadaten können in einem Graphen nie vollständig sein; sinnvoll ist relative Vollständigkeit. Modellbezug: Vollständigkeit kontextsensitiv operationalisieren. \\

K2.2 & Empfohlene Felder über Pflichtfelder hinaus
& Aussagen, die betonen, dass DCAT-AP.de-Pflichtfelder allein für gute Metadaten nicht ausreichen.
& Codieren, wenn empfohlene Felder, Soll-Kriterien oder optionale Angaben als qualitätssteigernd beschrieben werden.
& Die Möglichkeiten von DCAT-AP.de sollten möglichst umfassend genutzt werden. Modellbezug: empfohlene Felder im Score berücksichtigen. \\

K2.3 & Fachlich vorhandene Informationen übernehmen
& Aussagen, die fordern, dass bereits vorhandene fachliche Informationen nicht verworfen werden.
& Codieren, wenn vorhandene Ansprechpartner, fachliche Rollen, ergänzende Informationen oder domänenspezifische Angaben in Metadaten übernommen werden sollen.
& Fachlich vorhandene Informationen, etwa mehrere Ansprechpartner im Geobereich, sollten nicht entfernt werden. Modellbezug: Reichhaltigkeit ohne Informationsverlust. \\

K2.4 & Relevante Felder statt maximale Feldanzahl
& Aussagen, die eine gezielte Auswahl relevanter Felder gegenüber vollständiger Befüllung aller DCAT-Felder bevorzugen.
& Codieren, wenn nicht alle möglichen Felder, sondern die wichtigen Felder korrekt und sorgfältig ausgefüllt werden sollen.
& Nicht alle circa 70 DCAT-Felder sind für jeden Datensatz sinnvoll. Modellbezug: Gewichtung nach Relevanz statt bloßer Feldanzahl. \\

\midrule
\multicolumn{5}{l}{\textbf{K3 Formale und syntaktische Qualität}} \\
\midrule

K3.1 & DCAT-AP.de-Mindestvalidität
& Aussagen zur Erfüllung der Mindest- oder Pflichtanforderungen des DCAT-AP.de-Standards.
& Codieren, wenn Standardkonformität als Basiskriterium oder Mindestvoraussetzung beschrieben wird.
& DCAT-AP.de-Validität als Minimalvoraussetzung. Modellbezug: Basiskriterium formaler Qualität. \\

K3.2 & SHACL-Konformität als Basisdimension
& Aussagen zur Bedeutung von SHACL-Validierung für formale Qualitätsprüfung.
& Codieren, wenn SHACL als Prüfung der Standardkonformität, Fehlerfreiheit oder syntaktischen Korrektheit beschrieben wird.
& SHACL verbindet Feldvorhandensein und syntaktisch korrekte Ausfüllung. Modellbezug: formale Validierungsdimension. \\

K3.3 & URI- und Vokabularkorrektheit
& Aussagen zur korrekten Verwendung von URIs, kontrollierten Vokabularen und standardisierten Werten.
& Codieren, wenn richtige URIs, Vokabulare oder normierte Wertelisten als formale Qualitätsanforderung genannt werden.
& Korrekte URIs und Vokabulare als grundlegende Prüfung. Modellbezug: syntaktisch-formaler Indikator. \\

K3.4 & Trennung von fehlendem Feld und syntaktischem Fehler
& Aussagen, die unterschiedliche formale Fehlerarten voneinander trennen.
& Codieren, wenn gefordert wird, dass fehlende Felder und falsch ausgefüllte Felder getrennt ausgewiesen werden.
& Es soll klar sein, ob ein Problem syntaktisch ist oder ob ein Feld fehlt. Modellbezug: differenzierte Fehlerklassen. \\

K3.5 & XML/RDF-Validität und Datentypprüfung
& Aussagen zur technischen Validität von XML, RDF, Datentypen oder strukturellen Formaten.
& Codieren, wenn valides XML, RDF-Struktur, Datumsformate, Datentypen oder technische Formatvalidität thematisiert werden.
& Valides XML, korrekte URIs und Vokabulare als grundlegende Prüfungen. Modellbezug: technisch-syntaktische Prüflogik. \\

\midrule
\multicolumn{5}{l}{\textbf{K4 Semantische und inhaltliche Qualität}} \\
\midrule

K4.1 & Nichtssagende oder zu kurze Titel
& Aussagen zu Titeln, die vorhanden sind, aber den Datensatz nicht ausreichend beschreiben.
& Codieren, wenn Titel als zu kurz, generisch, mehrdeutig oder unverständlich beschrieben werden.
& „Los 1“ als Titel ohne erkennbare Aussage. Modellbezug: semantischer Titelindikator. \\

K4.2 & Titel-Beschreibung-Redundanz
& Aussagen, in denen die Beschreibung keinen eigenständigen Informationswert gegenüber dem Titel hat.
& Codieren, wenn Titel und Beschreibung identisch oder nahezu identisch sind.
& Titel und Beschreibung sind identisch; die Beschreibung liefert keinen Mehrwert. Modellbezug: Redundanzprüfung. \\

K4.3 & Regionale Begriffe und Fachsprache
& Aussagen zu sprachlichen Ausdrücken, die für Außenstehende schwer verständlich sind.
& Codieren, wenn regionale Begriffe, Fachsprache, domänenspezifische Terminologie oder fehlende allgemeinverständliche Synonyme genannt werden.
& „I-Dötzchen“ oder „eulitoral“ als schwer verständliche Begriffe. Modellbezug: Verständlichkeitsprüfung von Freitext. \\

K4.4 & Verständliche und erklärende Beschreibung
& Aussagen zur inhaltlichen Erklärungskraft von Beschreibungen.
& Codieren, wenn eine Beschreibung den Inhalt der Daten verständlich erklären und Nutzenden Orientierung geben soll.
& Beschreibung soll erklären, was die Daten enthalten und ob sie nutzbar sind. Modellbezug: semantische Beschreibungsqualität. \\

K4.5 & Inhaltliche Widersprüche zwischen Metadatenfeldern
& Aussagen zu unplausiblen oder widersprüchlichen Angaben innerhalb eines Metadatensatzes.
& Codieren, wenn Titel, Herausgeber, Ort, Beschreibung, Kategorie oder andere Metadatenfelder inhaltlich nicht zusammenpassen.
& Titel verweist auf Nürnberg, Herausgeber ist Köln. Modellbezug: Plausibilitätsprüfung zwischen Feldern. \\

K4.6 & Keywords, Themenfeld und Kategorie als semantische Prüffelder
& Aussagen, die Schlagwörter, Themenfelder oder Kategorien als relevante semantische Qualitätsmerkmale behandeln.
& Codieren, wenn Keywords, Kategorien, Themenfelder oder Schlagworte hinsichtlich Passung, Mehrwert oder Mapping bewertet werden.
& Keywords und Kategorien sollen den Inhalt passend beschreiben. Modellbezug: semantischer Indikator für thematische Passung. \\

\midrule
\multicolumn{5}{l}{\textbf{K5 Technische Zugänglichkeit, Aktualität und Nutzbarkeit}} \\
\midrule

K5.1 & Dead Links und Linkrot
& Aussagen zu nicht mehr erreichbaren Links oder veralteten Ressourcenverweisen.
& Codieren, wenn defekte URLs, Linkrot, nicht funktionierende Downloads oder nicht erreichbare Services genannt werden.
& Distributionslink, Download oder Datenservice funktioniert nicht mehr. Modellbezug: Verfügbarkeitsindikator. \\

K5.2 & Funktionierende Distributions-, Download- und Servicelinks
& Aussagen zur Erreichbarkeit und Nutzbarkeit von Zugriffspfaden.
& Codieren, wenn funktionierende Download-URLs, Access-URLs, Datenservices oder Distributionen als Qualitätsanforderung genannt werden.
& Nutzende sollen leicht zu den beschriebenen Datensätzen gelangen. Modellbezug: technische Zugänglichkeit. \\

K5.3 & Aktualität und Harvesting-Zeitpunkt
& Aussagen zu Alter, Aktualität, Harvesting-Datum oder Katalog-Metadaten.
& Codieren, wenn Aktualität, letzter Harvesting-Zeitpunkt, Catalog Records oder Pflegezustand als Qualitätsindikator beschrieben werden.
& Harvesting-Datum kann zeigen, dass ein Eintrag dem letzten Stand entspricht. Modellbezug: Aktualitäts- und Pflegeindikator. \\

K5.4 & Maschinenlesbarkeit und nicht-proprietäre Formate
& Aussagen zur technischen Weiterverarbeitbarkeit von Ressourcenformaten.
& Codieren, wenn maschinenlesbare, offene, nicht-proprietäre oder gut lesbare Formate als Nutzbarkeitsanforderung genannt werden.
& Bereitstellende fragen nach gut lesbaren und nicht-proprietären Dateiformaten. Modellbezug: Open-Data-Fitness. \\

\midrule
\multicolumn{5}{l}{\textbf{K6 Interoperabilität und Linked-Data-Qualität}} \\
\midrule

K6.1 & URI statt Literal
& Aussagen, die maschineninterpretierbare Identifikatoren gegenüber Freitextangaben bevorzugen.
& Codieren, wenn ein URI-, ID- oder Entitätsverweis als verlässlicher oder hochwertiger als ein Literal beschrieben wird.
& Eine GND-ID ist verlässlicher als ein Organisationsname als Freitext. Modellbezug: Linked-Data-Reifegrad. \\

K6.2 & Externe Referenzsysteme, Wikidata und GND
& Aussagen zu externen Normdaten und Wissenssystemen als Verknüpfungsressourcen.
& Codieren, wenn Wikidata, GND oder andere externe Referenzsysteme zur Anreicherung oder Identifikation genannt werden.
& Verknüpfungen zu Wikidata oder Gemeinsamer Normdatei. Modellbezug: externe semantische Referenzierung. \\

K6.3 & Kontrollierte Vokabulare und Thesauri
& Aussagen zur Verwendung kontrollierter Vokabulare, Thesauri oder standardisierter Begriffssysteme.
& Codieren, wenn Thesauri, EU-Vokabulare, UBA-Thesaurus oder kontrollierte Schlagwortsysteme genannt werden.
& Thesauri können Fachbegriffe beschreiben und Kontext erschließen. Modellbezug: Interoperabilität und semantische Harmonisierung. \\

K6.4 & Querverbindungen zwischen vergleichbaren Datensätzen
& Aussagen zu Verknüpfungen zwischen ähnlichen oder gleichen Datensätzen.
& Codieren, wenn gleiche, ähnliche oder vergleichbare Datensätze über Portale, Kommunen oder Domänen hinweg verknüpft werden sollen.
& Querverbindungen zwischen vergleichbaren Datensätzen aus Kommunen. Modellbezug: graphbasierte Vergleichbarkeit. \\

K6.5 & Organisationsgraph und Vermeidung von Blank Nodes
& Aussagen zur strukturierten Modellierung von Organisationen und zur Vermeidung schwacher RDF-Strukturen.
& Codieren, wenn Organisationsgraph, Agentenmodellierung, Blank Nodes oder strukturelle Linked-Data-Qualität genannt werden.
& Organisationen sollten sinnvoll über einen Organisationsgraphen abgebildet werden. Modellbezug: strukturelle RDF-Qualität. \\

\midrule
\multicolumn{5}{l}{\textbf{K7 Kontextsensitivität und Fit for Purpose}} \\
\midrule

K7.1 & Domänenspezifische Feldrelevanz
& Aussagen, nach denen bestimmte Metadatenfelder je nach Fachdomäne unterschiedlich relevant sind.
& Codieren, wenn Rollen, Relationen, Media Types, Geolocation oder andere Felder abhängig von Geodaten, Statistikdaten oder Fachdomänen bewertet werden.
& Creator und Maintainer sind bei Geodaten relevanter als bei Statistikdaten. Modellbezug: kontextsensitive Gewichtung. \\

K7.2 & Geodaten und Normierungsgrad
& Aussagen zum besonderen Normierungsgrad und zur spezifischen Qualitätslogik von Geodaten.
& Codieren, wenn Geodaten, ISO-Kontext, INSPIRE, GIS, Bounding Boxes oder normierte Fachstandards als besondere Domäne genannt werden.
& Geodaten sind stark normiert und erleichtern hochwertige Metadaten. Modellbezug: Domänenspezifik im Modell. \\

K7.3 & Föderale Vergleichbarkeit
& Aussagen zu Qualitätsvergleichen zwischen Ländern, Kommunen oder föderalen Datenbereitstellern.
& Codieren, wenn Schleswig-Holstein, Niedersachsen, Bayern, Länder- oder Kommunalvergleiche als Kontext für Bewertung genannt werden.
& Abweichung zwischen Küstenländern kann Qualitätsproblem sein, zwischen Küstenland und Bayern nicht zwingend. Modellbezug: föderaler Kontext. \\

K7.4 & Stakeholder-spezifische Qualitätsfragen
& Aussagen, nach denen unterschiedliche Akteure unterschiedliche Qualitätsinteressen haben.
& Codieren, wenn Qualitätsmetriken für Länder, Bund, Kommunen, Fachpersonen oder andere Stakeholder zugeschnitten werden sollen.
& Mindestens 17 relevante Stakeholder haben zugeschnittene Qualitätsfragen. Modellbezug: konfigurierbare Qualitätsansichten. \\

K7.5 & Unterschiedliche Perspektiven von Bereitstellenden und Nutzenden
& Aussagen, die Unterschiede zwischen Datenbereitstellenden und Datennutzenden betonen.
& Codieren, wenn dieselbe Metadatenqualität aus Bereitstellenden- und Nutzendenperspektive unterschiedlich bewertet wird.
& Ein alter Datensatz kann für Bereitstellende valide sein, für Nutzende aber verdächtig wirken. Modellbezug: perspektivenabhängige Interpretation. \\

\midrule
\multicolumn{5}{l}{\textbf{K8 Rohdatenbezug und Metadaten-Daten-Kongruenz}} \\
\midrule

K8.1 & Formatkongruenz zwischen Metadaten und Ressource
& Aussagen zum Abgleich zwischen angegebenem Format und tatsächlicher Ressource.
& Codieren, wenn geprüft werden soll, ob das in den Metadaten angegebene Format mit Datei, Service oder Inhalt übereinstimmt.
& Metadaten geben CSV an, aber der Inhalt ist kein CSV; oder CSV hängt an, während PDF angegeben ist. Modellbezug: Kongruenzindikator. \\

K8.2 & Serviceverhalten und tatsächliche Nutzbarkeit
& Aussagen zum Abgleich zwischen erwartetem und tatsächlichem Verhalten von Services.
& Codieren, wenn etwa WMS-/WFS-Dienste, APIs, XML-Fehlermeldungen oder erwartete Kartenansichten als Prüfgegenstand genannt werden.
& Mensch erwartet eine Karte, bekommt aber eine XML-Fehlermeldung. Modellbezug: praktische Ressourcenvalidierung. \\

K8.3 & Dateninhalt gegen Titel und Beschreibung prüfen
& Aussagen, nach denen tatsächliche Dateninhalte gegen Metadatenangaben geprüft werden.
& Codieren, wenn Titel, Beschreibung oder Themenangaben mit tatsächlichem Inhalt der Datei oder Ressource verglichen werden.
& Titel sagt „Vornamen“, Datensatz enthält hauptsächlich Zahlen. Modellbezug: semantischer Rohdatenabgleich. \\

K8.4 & Bounding Box und Flächendeckung aus Daten ableiten
& Aussagen zu geodatenbezogenen Informationen, die aus den Daten selbst berechnet oder geprüft werden können.
& Codieren, wenn Bounding Box, Flächendeckung oder räumliche Abdeckung aus Daten statt manueller Metadateneingabe abgeleitet werden sollen.
& Bounding Box ist menschlich eingetragen und fehleranfällig; besser wäre Berechnung aus Daten. Modellbezug: datenbasierte Kontextprüfung. \\

K8.5 & Fachliche Grenzen der Rohdatenbewertung
& Aussagen zu Grenzen einer Bewertung des eigentlichen Datensatzes durch ein Metadatenportal.
& Codieren, wenn fachliche Zuständigkeit, domänenspezifische Standards oder begrenzte Beratungskompetenz beim Rohdateninhalt thematisiert werden.
& Metadatenportal kann Rohdatenqualität nur eingeschränkt fachlich beurteilen. Modellbezug: Begrenzung des Prüfanspruchs. \\

\midrule
\multicolumn{5}{l}{\textbf{K9 Scoring, Gewichtung, Reporting und Verbesserungshinweise}} \\
\midrule

K9.1 & Wissenschaftlich begründete Metriken
& Aussagen, die eine theoretische, wissenschaftliche oder methodische Begründung von Qualitätsmetriken fordern.
& Codieren, wenn kritisiert wird, dass Scores Kriterien messen, ohne zu begründen, warum diese Kriterien relevant sind.
& Ein Score sollte belegen, warum genau diese Kriterien verwendet werden. Modellbezug: begründete Indikatorauswahl. \\

K9.2 & SHACL als Teildimension statt Gesamtscore
& Aussagen, die SHACL- oder DCAT-AP.de-Konformität als Teil, aber nicht als vollständiges Qualitätsurteil verstehen.
& Codieren, wenn formale Validierung als eigene Dimension im Gesamtscore beschrieben wird.
& DCAT-AP.de-Qualität sollte als eine Dimension im Gesamtscoring angezeigt werden. Modellbezug: separater formaler Teilschritt. \\

K9.3 & Teil-Scores zur Problemlokalisierung
& Aussagen zur Aufteilung eines Scores in Dimensionen oder Teilbewertungen.
& Codieren, wenn Teil-Scores, Kategorien oder getrennte Bewertungsbereiche zur Diagnose von Problemen gefordert werden.
& Teilbewertungen helfen zu erkennen, wo genau ein Problem liegt. Modellbezug: multidimensionaler Report. \\

K9.4 & Ranking als Verbesserungsanreiz
& Aussagen zur Nutzung von Rankings oder regelmäßigen Prüfungen zur Qualitätsverbesserung.
& Codieren, wenn Länder, Kommunen oder Portale sich vergleichen und daraus Verbesserungsbedarf ableiten sollen.
& Rankings können zeigen, wo Luft nach oben besteht. Modellbezug: Benchmarking, aber mit Vorsicht. \\

K9.5 & Verständliche Fehlermeldungen und Handlungshinweise
& Aussagen zur Verständlichkeit und praktischen Nutzbarkeit von Validierungs- oder Qualitätsreports.
& Codieren, wenn Fehler erklärt, konkret lokalisiert oder mit Verbesserungshinweisen versehen werden sollen.
& Eine reine Fehlermeldung ist wenig hilfreich; es braucht Erklärung und Behebungshinweis. Modellbezug: handlungsorientiertes Reporting. \\

\midrule
\multicolumn{5}{l}{\textbf{K10 KI-gestützte Analyseverfahren und Automatisierungsgrenzen}} \\
\midrule

K10.1 & Metadaten-Linter
& Aussagen zu KI- oder regelgestützten Systemen, die Metadaten prüfen und Hinweise geben.
& Codieren, wenn ein Linter, Prüfwerkzeug oder Vorschlagssystem auf Basis von DCAT-AP.de, SHACL oder Spezifikation beschrieben wird.
& Agentische KI als Metadaten-Linter. Modellbezug: unterstützende Qualitätsprüfung. \\

K10.2 & Freitextprüfung und Beschreibungsgüte
& Aussagen zur KI-gestützten Bewertung oder Verbesserung natürlicher Sprache in Metadaten.
& Codieren, wenn Titel, Beschreibungen, einfache Sprache, Synonyme, Fachsprache oder Style Guides mit KI-Unterstützung geprüft werden.
& KI kann Beschreibungstexte oder die Güte von Beschreibungstexten bewerten. Modellbezug: semantische Freitextanalyse. \\

K10.3 & Keyword-Mapping und Assisted Learning
& Aussagen zu KI- oder ML-gestützter Zuordnung von Stichworten zu Kategorien, Vokabularen oder Schlagwortsystemen.
& Codieren, wenn einfache Freitext-Stichworte auf Kategorien, echte Schlagworte oder EU-Vokabulare gemappt werden sollen.
& Assisted Learning kann qualitativ hochwertigere Keywords schaffen. Modellbezug: KI-gestützte Harmonisierung. \\

K10.4 & SPARQL- und Analyseunterstützung
& Aussagen zur KI-Unterstützung bei technischen Abfragen oder Analysebausteinen.
& Codieren, wenn KI zur Erstellung spezifischer SPARQL-Queries, Filter oder Analysebausteine eingesetzt werden soll.
& KI könnte Fachpersonen helfen, hochspezifische SPARQL-Queries zu entwickeln. Modellbezug: technische Assistenz statt Bewertungsautomat. \\

K10.5 & Semantischer Graph als bessere KI-Grundlage
& Aussagen, nach denen KI zuverlässiger arbeitet, wenn sie auf einem semantischen Graphen statt auf unstrukturierter Textbasis arbeitet.
& Codieren, wenn Graphdaten, semantische Verknüpfungen oder externe Graphen als Grundlage für KI-Nutzung beschrieben werden.
& Auf einem semantischen Graphen kann KI datenbasierter arbeiten und weniger halluzinieren. Modellbezug: graphbasierte Kontextgrundlage. \\

K10.6 & Rohdatenanalyse durch KI
& Aussagen zur KI-gestützten Analyse der tatsächlichen Ressourcen oder großer Datenmengen.
& Codieren, wenn KI in die Daten selbst hineinschauen, Muster erkennen, Strukturen vergleichen oder Rohdaten gegen Metadaten prüfen soll.
& KI kann große Datenmengen analysieren und Muster ableiten. Modellbezug: optionale kontextbezogene KI-Prüfung. \\

K10.7 & Datenhoheit bei automatisierter Verbesserung
& Aussagen zu Verantwortung, Schreibrechten oder Zuständigkeit bei automatisierter Veränderung von Metadaten.
& Codieren, wenn gefragt wird, wer Metadaten ändern, korrigieren oder anreichern darf.
& Wer hat die Hoheit über die Daten, wenn KI etwas verbessert oder verändert? Modellbezug: KI nur als Analyse- oder Empfehlungsebene. \\

K10.8 & Halluzination, stochastische Ausreißer und fachliche Fehlinterpretation
& Aussagen zu Unsicherheit, Fehleranfälligkeit, Halluzination oder fachlicher Fehlbewertung durch KI.
& Codieren, wenn stochastische Ausreißer, Halluzinationen, Mehrdeutigkeiten oder falsche Vereinfachungen durch KI thematisiert werden.
& KI kann „Watt“ nicht zuverlässig zwischen elektrischer Einheit und Wattenmeer unterscheiden. Modellbezug: keine autonome KI-Bewertungsinstanz. \\

K10.9 & Ressourcenaufwand und Betriebslast
& Aussagen zu Kosten, Rechenaufwand, Serverlast oder Betriebsvoraussetzungen KI-gestützter Verfahren.
& Codieren, wenn KI-Bewertungen wegen Wiederholungen, Last, Betrieb oder Ressourcenaufwand problematisiert werden.
& Wiederholte KI-Bewertungen können mehr Ressourcen kosten als sie Nutzen bringen. Modellbezug: Aufwand-Nutzen-Prüfung für KI-Komponenten. \\

\end{longtable}

\endgroup
